%%
%%Script
%%

\chapter{Scripting}
\label{ch:Scripting}

De eerste fase van het uitwerken van het project, begon met het maken van een script dat een gebruiker kon aanmaken in Active Directory. De basis van het script is een script dat je eenvoudig op internet kan vinden \autocite{Microsoft2024a}. Op vlak daarvan is het dan nodig om dit aan te passen naar de gebruikte ``Organizational Units'' binnen de Active Directory van de gemeente. \newline

De gemeente werkt met sites, op deze sites zijn soms onderverdelingen op vlak van diensten. Dit geldt echter niet altijd op deze manier. Zo is de site gemeentehuis bijvoorbeeld opgesplitst in diensten zoals personeelsdienst, vergunningen of informatica. Maar geldt deze structuur niet in de onthaalpoort, zij zitten aan de onthaalpoort aan het Donkmeer. Hier zitten de diensten van sport, jeugd, toerisme. Zij zitten wel allemaal samen in dezelfde OU. Dit komt omdat deze diensten nauwer samenwerken en vaak samen overleggen als 1 dienst.\newline 

Aangezien Power Automate, waar voor werd gekozen doordat een andere collega reeds aan de slag was hiermee in een gelijkaardig project, met een betalende licentie is, waren de initiële stappen op een andere manier aangepakt dan de uiteindelijke uitwerking. Eerst werd er dus gefocust op een script maken. Dit script haalde gegevens uit een Excel en zette die om naar een gebruiker in de Active Directory. \newline

De kolommen in Excel kregen de namen van de uiteindelijk gebruikte termen en variabelen die in de verdere fasen in Power Automate gebruikt worden. In deze kolommen kwamen waarden zoals voornaam, achternaam, startdatum, mogelijke einddatum, functietitel, site, dienst en het telefoonnummer.\newline 

Vanuit de waarden in deze kolommen werd er dus een gebruiker aangemaakt. Dit verliep allemaal via PowerShell. Dit script had de basis voor het uiteindelijk gebruikte script. De gebruikers van diensten krijgen standaard een groep mee. Deze groep geeft hen toegang tot bepaalde mappen op de fileserver. \newline

Ook wordt er via dit script een gebruikersnaam aangemaakt. Als je in de gemeente gaat kijken naar het maken van de gebruikersnaam, dan kom je op 2 verschillende gebruikersnamen. Zo heb je de ``User logon name'' en de ``User logon name (pre-Windows 2000)''. 

Bij de ``User logon name'' wordt er gebruik gemaakt van de email. Dit is dus bijvoorbeeld ``ewout.vanmullem@berlare.be''. Bij het aanmaken van een nieuwe gebruiker is dit standaard @intranet.berlare.be. Hier moest dus een aanpassing voor gemaakt worden.

Voor de ``User logon name (pre-Windows 2000)'', staat er standaard ``intranet\textbackslash'' met dan de voornaam.naam als gebruikersnaam. Op de gemeente staat de gebruikersnaam als achternaam + eerste letter van de voornaam. Volgens een auditor die een aantal jaar geleden langs is geweest, is dit veiliger aangezien dit moeilijker te achterhalen is dan voornaam + naam. De collega's gaan hier niet echt met akkoord maar aangezien vroegere collega's van de IT dienst dit al hadden doorgevoerd, werd er besloten hiermee gewoon verder te doen.  

\begin{figure}[htbp]
    \centering
    \includegraphics[width=0.4\textwidth]{OrganizationalUnit.png}
    \caption[OU structuur]{Deze afbeelding is de OU structuur die in de Active Directory van de gemeente gebruikt wordt. }
    \label{fig:OrganizationalUnit}
\end{figure}

%\begin{listing}[htbp]
%    \inputminted{powershell}{PowerAutomate.ps1}
%    \caption[PowerAutomate script]{Volledig PowerShell-script voor geautomatiseerde onboarding met power automate}
%    \label{lst:Power-Automate-script}
%\end{listing}